\documentclass{article}
\usepackage[utf8]{inputenc}
\usepackage{geometry}
\usepackage{datetime2}
\usepackage{xparse}
\usepackage{amsmath}
\usepackage{booktabs}
\usepackage{array}
\usepackage{hyperref}
\usepackage{listings}
\usepackage{xcolor}

 \geometry{
 a4paper,
 total={170mm,257mm},
 left=20mm,
 top=20mm,
 }
 
 \usepackage{graphicx}
 \usepackage{titling}

 \title {Tugas 3 - Fuzzy Inference System 
}

\date{16 Maret 2026}
 
 \usepackage{fancyhdr}
\fancypagestyle{plain}{%  the preset of fancyhdr 
    \fancyhf{} % clear all header and footer fields
    \fancyfoot[R]{}
    \fancyfoot[L]{}
    \fancyhead[L]{\includegraphics[width=2cm]{ITS_pojok.pdf}}
    \fancyhead[R]{ \DTMnow}
}
\makeatletter
\def\@maketitle{%
  \newpage
  \null
  \vskip 1em%
  \begin{center}%
  \let \footnote \thanks
    {\LARGE \@title \par}%
    \vskip 1em%
    %{\large \@date}%
  \end{center}%
  \par
  \vskip 1em}
\makeatother

\usepackage{lipsum}  
% \usepackage{cmbright} % Change the font by zee

% Pengaturan warna kode
\lstset{
    language=Python,
    basicstyle=\ttfamily,
    keywordstyle=\color{blue},
    commentstyle=\color{green},
    stringstyle=\color{red},
    numbers=left,
    numberstyle=\tiny,
    frame=single
}

\begin{document}

\maketitle

\noindent\begin{tabular}{@{}ll}
     Name &: Prayudhisti Basuki\\
     NRP &: 6022251075\\
     Teacher &: Muhammad Qomaruz Zaman, S.T., M.T., Ph.D.\\
     Class Name &: Kecerdasan Buatan (X)\\
     Github link &:  \href{https://github.com/kunthengx/tugas_kecerdasan_buatan/blob/c3d79eb01ee79165f0499d08e516e762eb2188eb/tugas3_fuzzy_inference.ipynb}{Tugas 3 - Fuzzy Inference System} \\
\end{tabular}
\vspace{0.75cm}

\section*{Deskripsi Tugas}

Rancang sebuah Fuzzy Inference System (FIS) untuk menentukan harga paket penjualan PV berdasarkan 

\begin{itemize}
    \item   Input 1 - Luasan Area PV : sedikit, sedang, banyak
    \item   Input 2 - Daya yang dihasilkan : kecil, menengah, besar
    \item   Output - Harga Paket: minimal 3 kelas (Basic, Premium, Platinum)
\end{itemize}

\setlength{\parindent}{0pt}
\textbf{Persyaratan !} Wajib ada aturan, misal :
\textbf{IF} PV sedikit \textbf{AND} Daya Besar \textbf{THEN} Paket = \textbf{Premium}\\
\\
Gunakan Model \textbf{Sugeno 0-Orde} (singleton output) dan \textbf{Defuzzyfikasi Weighted Average}
\\

\section*{PENYELESAIAN}
\subsection*{Perancangan Fuzzy Inference System (Sugeno Orde-0)}
Sistem ini bertujuan untuk menentukan harga paket berdasarkan dua input: \textbf{Luasan Area} dan \textbf{Daya}.
\\
\textbf{A. Fuzzifikasi (Variabel & Himpunan)}\\
Kita mendefinisikan semesta pembicaraan untuk setiap variabel:
\begin{enumerate}
    \item Input 1: \textbf{Luasan Area PV} (Skala 0-100 $m^2$)
    \begin{itemize}
        \item \textbf{Sedikit:} Kurva bahu kiri.
        \item \textbf{Sedang:} Kurva segitiga.
        \item \textbf{Banyak:} Kurva bahu kanan.
    \end{itemize}

    \item Input 2: \textbf{Daya yang Dihasilkan} (Skala 0-10 kWp)
    \begin{itemize}
        \item \textbf{Kecil:} Kurva bahu kiri.
        \item \textbf{Menengah:} Kurva segitiga.
        \item \textbf{Besar:} Kurva bahu kanan.
    \end{itemize}

    \item Output: \textbf{Harga Paket (Singleton)}
    \begin{itemize}
        \item \textbf{Basic:} 10 (Juta) Rupiah
        \item \textbf{Premium:} 25 (Juta) Rupiah
        \item \textbf{Platinum:} 45 (Juta) Rupiah
    \end{itemize}
\end{enumerate}

\textbf{B. Kurva Keanggotaan Fuzzy Inference System (Sugeno Orde-0)}\\
Dalam logika fuzzy, kurva himpunan keanggotaan ini sangat penting karena menentukan sejauh mana sebuah nilai input (seperti luas area) masuk ke dalam kategori tertentu (seperti "sedikit" atau "sedang"). \\

Fungsi keanggotaa yang digunakan adalah segitiga (\textit{triangular membership function}) yang didefinisikan oleh tiga parameter (\textit{a, b, c}) dimana $a \leq b \leq c$ serta nilai-nilai parameternya (\textit{a, b, c, d, d, e}) ditunjukkan sebagaimana gambar dibawah berikut

\begin{figure}[h]
    \centering
    \includegraphics[width=0.5\linewidth]{triangle membership function.png}
    %\caption{Fuzzyfikasi Suhu Ruangan}
    \label{fig:enter-label}
\end{figure}

        triangle z : a,b,c,d,e = 
        $\begin{cases}
        d, & x < a \\
        \frac{x - a}{b - a}, & a \leq x \leq b \\
        \frac{c - x}{c - b}, & b \leq x \leq c \\
        e, & x > c
        \end{cases}$
\\
Berikut adalah penjelasan dan visualisasi dari himpunan fuzzy yang kita gunakan:\\

\begin{figure}[h!]
    \centering
    \includegraphics[width=0.5\linewidth]{luas area.png}
    \caption{Himpunan Fuzzy Luasan Area PV}
    \label{fig:enter-label}
\end{figure}

\begin{figure}[h!]
    \centering
    \includegraphics[width=0.5\linewidth]{daya pv.png}
    \caption{Himpunan Fuzzy Daya PV yang Dihasilkan}
    \label{fig:enter-label}
\end{figure}

\begin{enumerate}
    \item Kurva Variabel Input (Luasan \& Daya)\\
Untuk input, kita menggunakan kombinasi kurva Trapesium di bagian ujung (karena nilai ekstrim biasanya memiliki derajat keanggotaan penuh 1) dan kurva Segitiga di bagian tengah.\\

\begin{itemize}
    \item \textbf{Kurva Segitiga} Bahu Kanan dan Kiri : Digunakan untuk kategori "Sedikit" dan "Banyak". Area datar di atas menunjukkan bahwa pada rentang tertentu, nilai tersebut 100\% dianggap "sedikit" atau "banyak".
    \item \textbf{Kurva Segitiga} : Digunakan untuk kategori "Sedang". Puncaknya menunjukkan nilai yang paling merepresentasikan kategori tersebut, sementara kemiringannya menunjukkan transisi ke kategori lain.
\end{itemize}

\begin{figure}[h!]
    \centering
    \includegraphics[width=0.5\linewidth]{harga paket.png}
    \caption{Himpunan Fuzzy Harga Paket (Singleton)}
    \label{fig:enter-label}
\end{figure}

    \item Kurva Output (Singleton Sugeno)\\
Pada model Sugeno Orde-0, output tidak berbentuk kurva lebar, melainkan berupa garis tunggal (Singleton) pada nilai konstanta tertentu.

\begin{itemize}
    \item \textbf{Basic:} Garis tegak lurus pada nilai 10.
    \item \textbf{Premium:} Garis tegak lurus pada nilai 25.
    \item \textbf{Platinum:} Garis tegak lurus pada nilai 45.
\end{itemize}
\end{enumerate}

\begin{lstlisting}[caption=Menghitung FIS Sugeno, language=Python]
def hitung_fis_sugeno(input_luas, input_daya):
    # Fuzzifikasi Input 1: Luas (0 - 100 m2)
    mu_luas = {
        'sedikit': triangle(input_luas, 25, 25, 50, 1, 0 ),
        'sedang' : triangle(input_luas, 25, 50, 75, 0, 0),
        'banyak' : triangle(input_luas, 50, 75, 100, 0, 1)
    }

    # Fuzzifikasi Input 2: Daya (0 - 10 kWp)
    mu_daya = {
        'kecil'   : triangle(input_daya, 2, 2, 5, 1, 0),
        'menengah': triangle(input_daya, 2, 5, 8, 0, 0),
        'besar'   : triangle(input_daya, 5, 8, 10, 0, 1)    
    }

    # Output Sugeno Orde-0 (Singleton dalam Juta Rp)
    Z = {'basic': 10, 'premium': 25, 'platinum': 45}

    # Aturan Fuzzy (Total 9 Aturan)
    # Format: (Luas, Daya, Output_Z)
    rules = [
        (mu_luas['sedikit'], mu_daya['kecil'],    Z['basic']),
        (mu_luas['sedikit'], mu_daya['menengah'], Z['basic']),
        (mu_luas['sedikit'], mu_daya['besar'],    Z['premium']),
        (mu_luas['sedang'],  mu_daya['kecil'],    Z['basic']),
        (mu_luas['sedang'],  mu_daya['menengah'], Z['premium']),
        (mu_luas['sedang'],  mu_daya['besar'],    Z['platinum']),
        (mu_luas['banyak'],  mu_daya['kecil'],    Z['premium']),
        (mu_luas['banyak'],  mu_daya['menengah'], Z['platinum']),
        (mu_luas['banyak'],  mu_daya['besar'],    Z['platinum'])
    ]

    # Inferensi & Defuzzifikasi (Weighted Average)
    numerator = 0
    denominator = 0

    for i, (m1, m2, z_val) in enumerate(rules):
        w = min(m1, m2)  # Operator AND (Mencari nilai minimum)
        if w > 0:
            numerator += w * z_val
            denominator += w

    if denominator == 0:
        return 0

    return numerator / denominator
\end{lstlisting}

\textbf{C. Aturan Fuzzy (Fuzzy Rules)}\\
Sesuai permintaan di soal, kita buat matriks aturan (total 9 aturan):\\

\begin{lstlisting}[caption=Membuat Rule Output, language=Python]
# Output Sugeno Orde-0 (Singleton dalam Juta Rp)
Z_print = {'basic': 10, 'premium': 25, 'platinum': 45}

rules_print = [
    (mu_luas_print['sedikit'], mu_daya_print['kecil'],    Z_print['basic']),
    (mu_luas_print['sedikit'], mu_daya_print['menengah'], Z_print['basic']),
    (mu_luas_print['sedikit'], mu_daya_print['besar'],    Z_print['premium']),
    (mu_luas_print['sedang'],  mu_daya_print['kecil'],    Z_print['basic']),
    (mu_luas_print['sedang'],  mu_daya_print['menengah'], Z_print['premium']),
    (mu_luas_print['sedang'],  mu_daya_print['besar'],    Z_print['platinum']),
    (mu_luas_print['banyak'],  mu_daya_print['kecil'],    Z_print['premium']),
    (mu_luas_print['banyak'],  mu_daya_print['menengah'], Z_print['platinum']),
    (mu_luas_print['banyak'],  mu_daya_print['besar'],    Z_print['platinum'])
]

for i, (m1, m2, z_val) in enumerate(rules_print):
    w = min(m1, m2)  # Operator AND (Mencari nilai minimum)
    if w > 0:
        print(f" - R{i+1} aktif dengan bobot (w): {w:.2f}, Nilai Z: {z_val}")
\end{lstlisting}

\begin{table}[h!]
\centering
\caption{Matriks Aturan Fuzzy}
\begin{tabular}{llll}
\hline
\textbf{Luasan / Daya} & \textbf{Kecil} & \textbf{Menengah} & \textbf{Besar}\\ \hline
\textbf{Sedikit}    & Basic     & Basic     & Premium     \\
\textbf{Sedang}     & Basic     & Premium   & Platinum     \\
\textbf{Banyak}     & Basic     & Platinum  & \textbf{Platinum}    \\ \hline
\end{tabular}
\label{tab:dummy_table} % Add a label for referencing
\end{table}

\textbf{D. Defuzzifikasi}\\
Karena memakai model Sugeno Orde-0, maka kita menggunakan \textbf{Weighted Average} : \\
\[
        Harga = \frac{\Sigma (w_{i} . z_{i})}{\Sigma w_{i}}
\]

Dimana $w_{i}$ adalah nilai keanggotaan (predikat) aturan ke-\textit{i}, dan $z_{i}$ adalah nilai konstanta output (singleton).

\begin{lstlisting}[caption=Menampilkan Hasil Perhitungan, language=Python]
# Get user input
user_luas = float(input("Masukkan Luas PV : "))
user_daya = float(input("Masukkan Daya PV : "))

harga_akhir = hitung_fis_sugeno(user_luas, user_daya)

print("-" * 30)
print(f"HASIL PERHITUNGAN SIMULASI")
print(f"Input Luas: {user_luas} m2 | Input Daya: {user_daya} kWp")
print(" ")
print("Detail Aktivasi Aturan:")

# Re-fuzzify inputs to get membership degrees for printing details
mu_luas_print = {
        'sedikit': triangle(user_luas, 25, 25, 50, 1, 0 ),
        'sedang' : triangle(user_luas, 25, 50, 75, 0, 0),
        'banyak' : triangle(user_luas, 50, 75, 100, 0, 1)
}

mu_daya_print = {
        'kecil'   : triangle(user_daya, 2, 2, 5, 1, 0),
        'menengah': triangle(user_daya, 2, 5, 8, 0, 0),
        'besar'   : triangle(user_daya, 5, 8, 10, 0, 1)
}
\end{lstlisting}



\end{document}
